% 《大学物理简明教程（第5版）》第1章：习题1.3.4—1.3.17。
% 扫描 PDF 实际第17—18页；书内印刷页17—18。
% 题意由扫描页人工核对，详解为独立推导；质量状态：待独立复核。

\sourceproblem{1.3.4}{17}{17}
\problemtext 已知平面运动的位置矢量
$\bm r(t)=2t\,\bm i+(4t^2-8)\,\bm j$（SI）。求轨迹、$t=1$至$2\,\mathrm s$的位移、$t=1\,\mathrm s$的速度与加速度（含大小、方向）。
\workings
由$x=2t,\ y=4t^2-8$消去$t$，得$y=x^2-8$，轨迹为抛物线。
\[
\Delta\bm r=\bm r(2)-\bm r(1)=(4\bm i+8\bm j)-(2\bm i-4\bm j)
=2\bm i+12\bm j\ {\rm m},\quad |\Delta\bm r|=2\sqrt{37}\ {\rm m}.
\]
\[
\bm v=\frac{d\bm r}{dt}=2\bm i+8t\bm j,\qquad
\bm a=\frac{d\bm v}{dt}=8\bm j.
\]
$t=1$时，$\bm v=2\bm i+8\bm j\ {\rm m\,s^{-1}}$，
$|\bm v|=2\sqrt{17}\ {\rm m\,s^{-1}}$，
方向与$x$正向夹角$\arctan4$；$\bm a=8\bm j\ {\rm m\,s^{-2}}$，沿$y$正向。
\crosscheck $\bm v$量纲为$\mathrm{m/s}$，$\bm a$量纲为$\mathrm{m/s^2}$；位置代回$y=x^2-8$一致（式中变量按SI数值表达）。

\sourceproblem{1.3.5}{17}{17}
\problemtext 质点在$Oxy$平面内满足$x=R\cos\omega t,\ y=R\sin\omega t$，$R,\omega>0$。求轨迹方程与任意时刻的位置、速度、加速度。
\workings
\[
x^2+y^2=R^2,\qquad
\bm r=R\cos\omega t\,\bm i+R\sin\omega t\,\bm j,
\]
\[
\bm v=R\omega(-\sin\omega t\,\bm i+\cos\omega t\,\bm j),
\qquad
\bm a=-R\omega^2(\cos\omega t\,\bm i+\sin\omega t\,\bm j)
=-\omega^2\bm r.
\]
因此为半径$R$、角速度$\omega$的匀速圆周运动，$v=R\omega,\ a=R\omega^2$。
\crosscheck $\bm r\cdot\bm v=0$，$\bm v\cdot\bm a=0$，加速度指向圆心。

\sourceproblem{1.3.6}{17}{17}
\problemtext 网球以$23.6\,\mathrm{m/s}$水平发出，击球点离地$2.37\,\mathrm m$，球网高$0.90\,\mathrm m$，水平距离$12\,\mathrm m$。
(1) 求球到网位置时高于网顶多少；(2) 若其他条件不变而发球方向向下偏$5.0^\circ$，判断能否过网。取$g=9.8\,\mathrm{m/s^2}$，忽略空气阻力。
\workings
(1) 水平初速度为$v_0$，达到网位置$t_1=12/v_0$，高度
\[
h_1=2.37-\frac{g}{2}\left(\frac{12}{23.6}\right)^2
\approx1.103\ {\rm m}.
\]
高于网顶$h_1-0.90\approx0.203\,\mathrm m$。
(2) 初始分量$v_{x0}=v_0\cos5^\circ,\ v_{y0}=-v_0\sin5^\circ$。
\[
t_2=\frac{12}{23.6\cos5^\circ},\qquad
h_2=2.37-12\tan5^\circ
-\frac{g}{2}\left(\frac{12}{23.6\cos5^\circ}\right)^2
\approx0.0436\ {\rm m}.
\]
因$h_2<0.90\,\mathrm m$，不能过网。
\crosscheck 两次到达球网平面之前球均未落地；结论不依赖球的半径或空气动力学，实际击球会受旋转和空气阻力影响。

\sourceproblem{1.3.7}{17}{17}
\problemtext 质点沿$x$轴运动，$x=4.5t^2-2t^3$（SI）。求第$2$秒内平均速度、$t=2\,\mathrm s$的瞬时速度、第$2$秒内路程。
\workings
$x(1)=2.5\,\mathrm m,\ x(2)=2.0\,\mathrm m$，故
\[
\bar v_{[1,2]}=\frac{x(2)-x(1)}{1\,\mathrm s}
=-0.5\,\mathrm{m/s}.
\]
$v(t)=9t-6t^2$，所以$v(2)=-6\,\mathrm{m/s}$。
在$(1,2)$内$v(t)=3t(3-2t)$于$t=1.5\,\mathrm s$改变符号。
$x(1.5)=3.375\,\mathrm m$，故路程
\[
s=[x(1.5)-x(1)]+[x(1.5)-x(2)]
=0.875+1.375=2.25\ {\rm m}.
\]
\crosscheck 平均速度可为负，而路程恒非负；不能以$|x(2)-x(1)|=0.5\,\mathrm m$代替路程。

\sourceproblem{1.3.8}{17--18}{17--18}
\problemtext 两车从同一起点同方向同时出发，其位置（SI）为
$x_A=4t+t^2$与$x_B=2t^2+2t^3$。判断刚出发谁在前方，求以后两车位置再次相同的时刻，以及$B$相对$A$的速度为零的时刻。
\workings
\[
x_A-x_B=t(4-t-2t^2).
\]
刚出发$t>0$足够小时该差为正，所以$A$在前。
除去$t=0$，两车再次相遇满足$2t^2+t-4=0$，故
\[
t_{\rm meet}=\frac{\sqrt{33}-1}{4}\ {\rm s}\approx1.186\ {\rm s}.
\]
速度为$v_A=4+2t,\ v_B=4t+6t^2$，
令$v_B-v_A=6t^2+2t-4=0$，正时间解为
\[
t_{\rm rel,0}=\frac23\ {\rm s}.
\]
\crosscheck $\mathrm d(x_B-x_A)/\mathrm dt=v_B-v_A$；相对速度为零是相对位移取得极值的时刻，不等于再次相遇。

\sourceproblem{1.3.9}{18}{18}
\problemtext 质点在水平圆轨道上运动，$R=2\,\mathrm m$，角速度满足$\omega=kt^2$。已知$t=2\,\mathrm s$时速率为$32\,\mathrm{m/s}$，求$t=1\,\mathrm s$的速率和加速度大小。
\workings
由$v=R\omega=Rkt^2$，有$32=2k\cdot4$，故$k=4\,\mathrm{rad/s^3}$。
$t=1$时，$\omega=4\,\mathrm{rad/s}$，$v=R\omega=8\,\mathrm{m/s}$；
角加速度$\alpha=\mathrm d\omega/\mathrm dt=2kt=8\,\mathrm{rad/s^2}$。
\[
a_\tau=R\alpha=16\ {\rm m/s^2},\quad
a_n=R\omega^2=32\ {\rm m/s^2},\quad
a=\sqrt{a_\tau^2+a_n^2}=16\sqrt5\ {\rm m/s^2}.
\]
\crosscheck $a_\tau$与$a_n$垂直，合成必须取平方和的平方根。

\sourceproblem{1.3.10}{18}{18}
\problemtext 石头在空气中由静止下落，沿下落正方向的加速度满足$a=A-Bv$，$A,B$均为正常量。求$v(t)$。
\workings
令$t=0$时$v=0$。微分方程$\mathrm dv/\mathrm dt+Bv=A$的积分因子为$e^{Bt}$，
\[
\frac{\mathrm d}{\mathrm dt}(ve^{Bt})=Ae^{Bt}
\quad\Longrightarrow\quad
\boxed{v(t)=\frac AB(1-e^{-Bt})}.
\]
\crosscheck $v(0)=0$，$v'(0)=A$；$t\to\infty$时$v\to A/B$，且$A-Bv\to0$，符合终端速度物理图像。

\sourceproblem{1.3.11}{18}{18}
\problemtext 一质点沿$x$轴运动，$a(x)=2+6x^2$（SI）；在$x=0$处速率为$10\,\mathrm{m/s}$，求运动到任一可达坐标处的速率。
\workings
利用$a=v\,\mathrm dv/\mathrm dx$，有
\[
v\,\mathrm dv=(2+6x^2)\,\mathrm dx,\qquad
\frac12v^2=2x+2x^3+C.
\]
由$v(0)^2=100$得$C=50$，故
\[
\boxed{|v(x)|=\sqrt{100+4x+4x^3}\ {\rm m/s}}
\]
（$x$按米的数值代入，且根号内非负）。
\crosscheck $a(x)>0$，因此速度的代数值随时间严格增加；负向初速度和转向情况须分段选取$v=\pm|v|$，不能仅凭题面给的初始速率推断全程运动方向。

\sourceproblem{1.3.12}{18}{18}
\problemtext 沿直线运动，$a(t)=4+3t$（SI），初始$x_0=5\,\mathrm m,\ v_0=0$，求$t=10\,\mathrm s$的速度与位置。
\workings
逐次积分并施加初值，
\[
v(t)=4t+\frac32t^2,\qquad
x(t)=5+2t^2+\frac12t^3.
\]
所以
\[
\boxed{v(10)=190\,\mathrm{m/s},\qquad x(10)=705\,\mathrm m.}
\]
\crosscheck $x'(t)=v(t)$且$v'(t)=4+3t$。

\sourceproblem{1.3.13}{18}{18}
\problemtext 直线运动质点的加速度$a=-A\omega^2\cos\omega t$，其中$A,\omega>0$；$t=0$时$v=0,\ x=A$。求运动方程。
\workings
\[
v(t)=\int a\,dt=-A\omega\sin\omega t+C_1,\quad C_1=0.
\]
再积分得$x(t)=A\cos\omega t+C_2$，由$x(0)=A$得$C_2=0$：
\[
\boxed{x(t)=A\cos\omega t.}
\]
\crosscheck $x''(t)=-A\omega^2\cos\omega t$且两个初值均成立。

\sourceproblem{1.3.14}{18}{18}
\problemtext 质点沿$R=1\,\mathrm m$圆周运动，角位置$\theta=2+3t^3$（rad，$t$单位s）。
求$t=2\,\mathrm s$的切向、法向加速度；以及加速度方向与半径所在直线成$45^\circ$时的角位移。
\workings
\[
\omega=\dot\theta=9t^2,\quad \alpha=\dot\omega=18t,\quad
a_\tau=R\alpha=18t,\quad a_n=R\omega^2=81t^4.
\]
$t=2$时$a_\tau=36\,\mathrm{m/s^2}$，$a_n=1296\,\mathrm{m/s^2}$。
与半径\emph{直线}夹角$45^\circ$要求$|a_\tau|=a_n$，对$t>0$有
$18t=81t^4$，故$t^3=2/9$，此时角位移
\[
\boxed{\Delta\theta=\theta(t)-\theta(0)=3t^3=\frac23\ {\rm rad}.}
\]
\crosscheck 指向圆心的法向加速度与向外径矢成钝角，故题面“与半径成$45^\circ$”按与\emph{半径所在直线}的锐角理解；若明确要求与\emph{向外半径矢量}成$45^\circ$，则无这样的时刻。

\sourceproblem{1.3.15}{18}{18}
\problemtext 质点在半径$R$的圆周上运动，弧长$s(t)=v_0t-\frac12bt^2$，$v_0,b$为常量。求$t$时的加速度；并求何时法向加速度的\emph{数值}等于$b$。
\workings
在$v_0>0,b>0$且弧长随时间增加的阶段$0\le t\le v_0/b$，
\[
v=\dot s=v_0-bt,\qquad
a_\tau=\dot v=-b,\qquad a_n=\frac{(v_0-bt)^2}{R}.
\]
加速度大小为
\[
a=\sqrt{b^2+\frac{(v_0-bt)^4}{R^2}}.
\]
由$a_n=b$且$v_0-bt\ge0$，
\[
\boxed{t=\frac{v_0-\sqrt{bR}}{b}}
\]
仅当$v_0\ge\sqrt{bR}$才有上述时间区间内的解。
\crosscheck 若$v_0<\sqrt{bR}$，减速至停下之前$a_n<b$，无非负解；不能未经有效域判断直接采用平方根的两支。

\sourceproblem{1.3.16}{18}{18}
\problemtext 半径$R=0.4\,\mathrm m$的圆轨道上，质点从静止开始以恒定角加速度$\alpha=0.2\,\mathrm{rad/s^2}$转动。求$t=2\,\mathrm s$的切向、法向及总加速度。
\workings
$\omega=\alpha t=0.4\,\mathrm{rad/s}$，
\[
a_\tau=R\alpha=0.08\ {\rm m/s^2},\qquad
a_n=R\omega^2=0.064\ {\rm m/s^2},
\]
\[
\boxed{a=\sqrt{0.08^2+0.064^2}\approx0.10245\ {\rm m/s^2}.}
\]
\crosscheck 切向与法向互相垂直，故不能简单将二者数值相加。

\sourceproblem{1.3.17}{18}{18}
\problemtext 一船以$30\,\mathrm{km/h}$向东航行，一小艇以$40\,\mathrm{km/h}$向北航行。求船上观察小艇与小艇上观察船的相对速度大小。
\workings
设东、北分别为$\bm i,\bm j$，
$\bm v_{\rm 船}=30\bm i,\ \bm v_{\rm 艇}=40\bm j$（$\mathrm{km/h}$）。
\[
\bm v_{\rm 艇/船}=40\bm j-30\bm i,\qquad
\bm v_{\rm 船/艇}=-\bm v_{\rm 艇/船}.
\]
因此两者观察到的相对速率相同：
\[
\boxed{|\bm v_{\rm 艇/船}|=|\bm v_{\rm 船/艇}|
=\sqrt{30^2+40^2}=50\ {\rm km/h}.}
\]
\crosscheck 两个相对速度矢量方向相反、模相等，满足伽利略相对速度关系。
